\chapter{Συσταδοποίηση}
\label{chap:clustering}

Η συσταδοποίηση εφαρμόζεται σε δύο επίπεδα, σύμφωνα με τη δομή που προτάθηκε
από την επίβλεψη: μια πρώτη ομαδοποίηση ως προς τις παραμέτρους με τη
μεγαλύτερη σημασία, δηλαδή την περιοχή και τον χρόνο, και εντός κάθε συστάδας
μια δεύτερη ομαδοποίηση ως προς επιπρόσθετες παραμέτρους. Το αποτέλεσμα είναι
δένδρο συστάδων, από κάθε κόμβο του οποίου μπορεί να εξαχθεί ανεξάρτητη
πληροφορία.

\section{Χαρακτηριστικά και κανονικοποίηση}

Στο πρώτο επίπεδο κάθε περιοχή περιγράφεται από τα έσοδα, το πλήθος
παραγγελιών, το πλήθος πελατών και το πλήθος των μηνών με δραστηριότητα. Το
τελευταίο χαρακτηριστικό είναι εκείνο που καθιστά το σύνολο χωροχρονικό και
όχι αμιγώς χωρικό: δύο περιοχές με ίδια έσοδα αποτελούν ουσιωδώς διαφορετικές
αγορές όταν η μία τα πραγματοποίησε σταθερά μέσα στο έτος και η άλλη
συγκεντρωμένα σε έναν μήνα.

Στο δεύτερο επίπεδο κάθε πελάτης περιγράφεται από τα χαρακτηριστικά RFM που
ορίστηκαν στην ενότητα~\ref{sec:measured-economics}.

Πριν από κάθε πέρασμα τα χαρακτηριστικά τυποποιούνται. Χωρίς τυποποίηση, ένα
χαρακτηριστικό μετρημένο σε νομισματικές μονάδες θα κυριαρχούσε έναντι ενός
μετρημένου σε πλήθη αποκλειστικά λόγω της κλίμακάς του, και ο αλγόριθμος θα
ομαδοποιούσε ως προς τις μονάδες μέτρησης αντί ως προς τη συμπεριφορά. Η
ιδιότητα αυτή επαληθεύεται με έλεγχο στον οποίο η πραγματική δομή βρίσκεται
εξ ολοκλήρου στο χαρακτηριστικό μικρής κλίμακας.

\section{Επιλογή του πλήθους συστάδων}

Το πλήθος των συστάδων δεν καθορίζεται εκ των προτέρων. Για κάθε υποψήφια τιμή
$k$ υπολογίζεται ο δείκτης περιγράμματος \parencite{rousseeuw1987silhouettes}
και επιλέγεται η τιμή με την υψηλότερη βαθμολογία. Η κατάτμηση προκύπτει
επομένως από τα δεδομένα και όχι από παραδοχή, ενώ οι βαθμολογίες όλων των
υποψηφίων τιμών καταγράφονται και παρατίθενται.

Όταν το πλήθος των εγγραφών δεν επαρκεί για ουσιαστική κατάτμηση, οι εγγραφές
αναφέρονται ως μία συστάδα χωρίς δείκτη περιγράμματος. Η επιλογή αυτή είναι
συνειδητή: προτιμάται η δήλωση της αδυναμίας από την επινόηση δομής που τα
δεδομένα δεν στηρίζουν.

\section{Πρώτο επίπεδο: περιοχή και χρόνος}

% TODO(PR8): Αποτελέσματα και σχολιασμός μετά την εισαγωγή του πλήρους
% συνόλου δεδομένων.

\section{Δεύτερο επίπεδο: υποσυστάδες RFM}

Εντός κάθε συστάδας πρώτου επιπέδου συγκεντρώνονται οι πελάτες όλων των
περιοχών που την απαρτίζουν και ομαδοποιούνται εκ νέου. Η ομαδοποίηση γίνεται
επί του συνόλου της συστάδας και όχι ανά περιοχή χωριστά, ώστε το δεύτερο
επίπεδο να περιγράφει τη συμπεριφορά της αγοράς ως ενιαίας οντότητας.

Οι πελάτες χωρίς καταγεγραμμένη περιοχή δεν εντάσσονται στο δένδρο, καθώς δεν
αντιστοιχούν σε κανέναν κλάδο του. Αναφέρονται μέσω της επίπεδης ανάλυσης RFM,
αντί να αποδοθούν αυθαίρετα σε κάποια ομάδα.

\section{Ερμηνεία των συστάδων}

Κάθε κόμβος συνοδεύεται από τους μέσους όρους των χαρακτηριστικών του, οι
οποίοι αποτελούν και τη βάση της ερμηνείας του. Οι συστάδες ταξινομούνται κατά
φθίνον μέγεθος, ώστε η αρίθμηση να παραμένει σταθερή και συγκρίσιμη μεταξύ
εκτελέσεων. Ο σπόρος τυχαιότητας είναι σταθερός, επομένως τα αποτελέσματα που
παρατίθενται είναι πλήρως αναπαραγώγιμα.
